% Abstract — mesmo formato do resumo (Guia UECE p. 47-49), com "Keywords".
% Tradução da primeira versão do resumo (18/11/2026); manter os dois em
% sincronia. Números: os mesmos do resumo (ponto decimal em inglês).
\begin{resumo}[ABSTRACT]
  \begin{otherlanguage*}{english}
    Agents based on language models generate SQL queries from a textual
    description of the database, and the correctness of these queries
    depends on what that description contains. This work measures the
    effect of three additive levels of context on the correctness of the
    generated SQL: the technical definition of the tables; the same
    definition with catalog descriptions of tables and columns; and the
    catalog extended with a textual semantic layer containing grain rules,
    temporal rules, relationship rules, synonyms and response style. The
    experiment uses a real, de-identified commercial database with 11
    PostgreSQL tables, 40 questions in Portuguese and two open-weight
    models from different families, \textit{gpt-oss-120b} and
    \textit{qwen3.8-27b}, accessed through a free API. Each question is
    answered under the three conditions, correctness is checked by
    comparing the result of the generated query with that of a reference
    query, and the conditions are compared with the exact McNemar test
    with Holm adjustment. With \textit{qwen3.8-27b}, accuracy goes from
    57.5\% to 60.0\% and to 87.5\%, and the semantic condition outperforms
    the catalog condition (adjusted \textit{p} of 0.002) and the technical
    one (adjusted \textit{p} of 0.001). With \textit{gpt-oss-120b},
    accuracy goes from 62.5\% to 70.0\% and to 82.5\%, with no significant
    difference after adjustment (smallest adjusted \textit{p} of 0.064).
    The catalog does not differ from the technical condition in either
    model. The errors fixed only by the semantic layer involve the
    handling of orphan references and the return of columns that were not
    requested, and the median latency of the semantic condition, measured
    on the provider's free tier, is about 5 times that of the technical
    one. The conclusion is that correctness
    is raised by rules on how the database should be queried, not by
    column descriptions; that a text file attached to the prompt
    reproduces the pattern reported by industry with proprietary
    platforms, with no license cost; and that this context acts as an
    instruction, not as a guarantee.

    \vspace{\onelineskip}
    \noindent\textbf{Keywords}: text-to-SQL; language models; semantic
    layer; prompt engineering; data catalog.
  \end{otherlanguage*}
\end{resumo}
